\documentclass[11pt,a4paper]{article}

\usepackage[utf8]{inputenc}
\usepackage[T1]{fontenc}
\usepackage[french]{babel}
\usepackage{lmodern}
\usepackage{amsmath,amssymb}
\usepackage{bm}
\usepackage{booktabs,array}
\usepackage[margin=2.5cm]{geometry}
\usepackage{hyperref}

\newcommand{\vm}{\bm{m}}
\newcommand{\vp}{\bm{p}}
\newcommand{\vs}{\bm{s}}
\newcommand{\vv}{\bm{v}}
\newcommand{\vq}{\bm{q}}
\newcommand{\mG}{\mathsf{G}}
\newcommand{\mA}{\mathsf{A}}
\newcommand{\Id}{\mathsf{I}}

\title{Polarisation du courant de spin dans le modèle de Slonczewski}
\author{}
\date{\today}

\begin{document}

\maketitle

\part{Couche libre mince : le modèle de Slonczewski}

\section{Le cadre}

On considère une tricouche F1 / espaceur non magnétique / F2 :
\begin{itemize}
  \item F1 : couche de référence, d'aimantation fixe $\vp$ (le
    \og polariseur \fg) ;
  \item F2 : couche libre, d'aimantation $\vm$ ;
  \item $\theta$ : angle entre $\vp$ et $\vm$, avec $\cos\theta = \vm\cdot\vp$.
\end{itemize}
Un courant de charge traverse l'empilement. Les électrons sortent de F1
avec une polarisation de spin, traversent l'espaceur, puis arrivent sur F2.

\section{Ce qui ne dépend pas de l'aimantation : \texorpdfstring{$P$}{P}}

$P$ est une propriété du matériau et de l'interface. Il vient de
l'asymétrie entre les électrons majoritaires et minoritaires : densité
d'états au niveau de Fermi, transmission dépendante du spin.

Par contre, $P$ n'a de sens que par rapport à un axe. Dans un ferro,
l'axe de quantification du spin est l'aimantation locale. Dire
\og le courant est polarisé à $P$\,\% \fg{} sous-entend donc
\og le long de $\vm$ à cet endroit \fg.

\section{Pourquoi c'est l'aimantation \emph{à l'interface} qui compte}

Quand un électron polarisé selon $\vp$ arrive sur F2, on décompose son
spin dans la base de $\vm$ :
\begin{itemize}
  \item \textbf{La composante parallèle à $\vm$} est majoritaire ou
    minoritaire dans F2. Elle est transmise ou réfléchie selon les
    coefficients de F2, et elle se propage sans problème.
  \item \textbf{La composante transverse à $\vm$} n'est pas un état propre
    dans F2. Elle disparaît en quelques couches atomiques, à cause de trois
    mécanismes :
    \begin{enumerate}
      \item le filtrage en spin à l'interface (réflexion et transmission
        différentes) ;
      \item la précession autour du champ d'échange, à des vitesses
        différentes selon les électrons ;
      \item le déphasage quand on moyenne sur la surface de Fermi, qui fait
        que ces précessions se compensent.
    \end{enumerate}
\end{itemize}

Ce moment angulaire transverse perdu est transféré à l'aimantation :
c'est le couple de transfert de spin, de forme
\begin{equation}
  \bm{\tau} \propto \vm\times(\vm\times\vp).
\end{equation}

Comme tout se passe sur environ $1\,\mathrm{nm}$, c'est \textbf{un effet
d'interface}. Seule l'orientation de $\vm$ juste à l'interface intervient.
C'est le sens de \og la polarisation dépend de l'aimantation à
l'interface \fg{} : la décomposition du spin entre composante absorbée et
composante transmise se fait sur l'aimantation locale de l'interface, pas
sur l'aimantation moyenne de la couche.

\section{Pourquoi l'efficacité dépend de \texorpdfstring{$\theta$}{θ}}

On pourrait croire que le courant de spin dans l'espaceur vaut simplement
$P\,J/e$ selon $\vp$. \textbf{C'est faux}, à cause des \textbf{réflexions
multiples} :
\begin{itemize}
  \item les électrons réfléchis par F2 reviennent avec une polarisation
    liée à $\vm$ ;
  \item ils sont à nouveau filtrés par F1, et ainsi de suite.
\end{itemize}

Le courant de spin qui arrive à l'interface de F2 dépend donc à la fois de
$\vp$ et de $\vm$. Slonczewski (1996) en tire une polarisation effective
qui dépend de l'angle :
\begin{equation}
  g(\theta) = \left[-4 + \frac{(1+P)^3\,(3+\cos\theta)}{4\,P^{3/2}}\right]^{-1}.
  \label{eq:g}
\end{equation}
Cette fonction est \textbf{asymétrique} : $g(\pi) > g(0)$. Il faut moins
de courant pour basculer depuis l'état antiparallèle que depuis l'état
parallèle.

Dans les codes micromagnétiques (mumax3, OOMMF), on utilise souvent la
forme généralisée :
\begin{equation}
  \varepsilon(\theta) = \frac{P\,\Lambda^2}
  {(\Lambda^2+1) + (\Lambda^2-1)\,(\vm\cdot\vp)}.
  \label{eq:eps}
\end{equation}
Avec $\Lambda = 1$, l'asymétrie disparaît et $\varepsilon = P/2$.

\section{Le couple final}

\begin{equation}
  \bm{\tau} = \frac{\gamma\hbar J}{2e\,\mu_0 M_s\,d}\;
  \varepsilon(\theta)\;\vm\times(\vm\times\vp).
  \label{eq:tau}
\end{equation}

Le facteur $1/d$ ($d$ = épaisseur de la couche libre) traduit le fait
qu'un couple déposé à l'interface est partagé sur toute la couche. Cela
suppose $\vm$ uniforme dans l'épaisseur (approximation macrospin dans
l'épaisseur). C'est pour ça qu'on confond \og $\vm$ à l'interface \fg{}
et \og $\vm$ de la couche \fg.

\section*{En une phrase}

\begin{quote}
  $P$ est fixée par le matériau. \textbf{L'axe} selon lequel on la mesure,
  et donc ce qui est absorbé, est fixé par $\vm$ à l'interface.
  \textbf{L'amplitude effective} dépend de l'angle entre $\vm$ et $\vp$, à
  cause des réflexions multiples entre les deux couches.
\end{quote}

\part{Couche libre épaisse : LLG et diffusion de spin}

\section{Le modèle de départ et ses limites}

On suppose maintenant F2 épaisse. L'aimantation $\vm(\bm{r},t)$ suit
l'équation LLG et l'accumulation de spin $\vs(\bm{r},t)$ suit l'équation
de diffusion de spin de Zhang--Levy--Fert (ZLF) :
\begin{equation}
  \partial_t \vs + \nabla\cdot Q
  = -\frac{\vs}{\tau_{sf}} - \frac{\vs\times\vm}{\tau_J}.
  \label{eq:zlf}
\end{equation}
On écrit le courant de spin comme un tenseur $Q$ de composantes
$Q_{i\alpha}$, avec $i$ l'indice d'espace (direction de l'écoulement) et
$\alpha$ l'indice de spin. Les indices latins désignent l'espace, les
indices grecs le spin, et on somme sur les indices répétés :
\begin{equation}
  Q_{i\alpha} = -\beta\frac{\mu_B}{|e|}\,j_{e,i}\,m_\alpha
  - 2D_0\left[\partial_i s_\alpha - \beta\beta'\,m_\alpha\,(\vm\cdot\partial_i\vs)\right],
  \qquad (\nabla\cdot Q)_\alpha = \partial_i Q_{i\alpha},
  \label{eq:courant}
\end{equation}
où $\bm{j}_e = \sigma\bm{E}$ est la densité de courant de charge. Le signe
moins vient de ce que le spin est transporté par les électrons, qui se
déplacent en sens inverse de $\bm{j}_e$. Sur les interfaces, on utilise plus
loin le scalaire $J$, positif pour des électrons qui vont de F1 vers F2 :
c'est le flux d'électrons, $J = -\bm{j}_e\cdot\hat{\bm{x}}$, l'axe
$\hat{\bm{x}}$ étant orienté de F1 vers F2. Le couple sur l'aimantation est $\bm{\tau}\propto\vm\times\vs$. Dans
FeeLLGood, $\beta=\beta'=P$ : la diffusion longitudinale vaut
$2D_0(1-P^2)$ et la diffusion transverse $2D_0$ (voir
section~\ref{sec:F1-feellgood}).

On note $\Gamma$ l'interface espaceur/F2, $\bm{n}$ la normale sortante de
F2 et $\vq = -(\bm{n}\cdot Q)|_\Gamma$, de composantes $q_\alpha = -n_iQ_{i\alpha}$,
le courant de spin qui
\emph{entre} dans F2. Le modèle de départ impose un flux (condition de
Neumann) :
\begin{equation}
  \vq = P\,\frac{\mu_B}{e}\,J\,\vp.
  \label{eq:neumann}
\end{equation}
Ce flux ne dépend ni de $\vm_0 = \vm|_\Gamma$, ni de $\vs|_\Gamma$. Le
modèle oublie donc trois choses :
\begin{itemize}
  \item \textbf{le filtrage par l'interface F2} : la partie du courant de
    spin portée par le courant de charge est polarisée selon $\vm_0$, pas
    selon $\vp$ ;
  \item \textbf{le retour de spin} (\emph{backflow}) : l'accumulation de
    spin qui se forme dans F2 renvoie du spin vers l'espaceur, ce qui
    modifie le flux injecté ;
  \item \textbf{les réflexions multiples entre F1 et F2}, qui donnent la
    dépendance en angle de type $g(\theta)$ ou $\varepsilon(\theta)$.
\end{itemize}

\section{Condition d'interface issue de la théorie des circuits}

\subsection{Conductances d'interface}

On décrit une interface N/F d'aimantation $\bm{u}$ par :
\begin{itemize}
  \item $G_\uparrow$, $G_\downarrow$ : conductances de spin par unité de surface ;
  \item $G_{\uparrow\downarrow} = G_r + iG_i$ : conductance de mélange
    (\emph{spin-mixing conductance}) ;
  \item $P_I = (G_\uparrow-G_\downarrow)/(G_\uparrow+G_\downarrow)$ :
    polarisation d'interface.
\end{itemize}
On les regroupe dans une matrice $3\times3$ qui dépend de $\bm{u}$ :
\begin{equation}
  \mG(\bm{u}) = G_\parallel\,\bm{u}\bm{u}^{\mathsf{T}}
  + G_r\left(\Id-\bm{u}\bm{u}^{\mathsf{T}}\right)
  - G_i\,[\bm{u}]_\times,
  \qquad [\bm{u}]_\times\,\bm{w} = \bm{u}\times\bm{w},
  \label{eq:G}
\end{equation}
avec $G_\parallel\propto(1-P_I^2)(G_\uparrow+G_\downarrow)$. Les trois
termes correspondent, dans l'ordre, à la diffusion longitudinale, à
l'absorption transverse et à la précession dans le champ d'échange
d'interface. Les conductances sont ici converties dans les unités de $\vs$,
et les facteurs de convention y sont absorbés.

\subsection{Courant de spin à travers l'interface}

On note $\vs_N$ l'accumulation de spin dans l'espaceur et
$\vs_F = \vs|_\Gamma$ celle dans F2 au bord. On élimine la chute de
potentiel de charge à l'aide de $J$ et on obtient :
\begin{equation}
  \boxed{\;
  \vq = P_I\,\frac{\mu_B}{e}\,J\,\vm_0
  + \mG(\vm_0)\,\left(\vs_N-\vs_F\right)
  \;}
  \label{eq:robin}
\end{equation}
Trois remarques :
\begin{itemize}
  \item La polarisation portée par le courant de charge est \textbf{selon
    $\vm_0$}. L'information sur $\vp$ n'arrive que par $\vs_N$,
    l'accumulation créée par F1 dans l'espaceur.
  \item La condition \eqref{eq:robin} est \textbf{linéaire en $\vs$} : c'est
    une condition de Robin. Ses coefficients dépendent de $\vm_0$.
  \item La théorie des circuits suppose un ferro fort : la partie
    transverse de $\vs_F$ est nulle à l'interface, et le courant transverse
    ne dépend que de $\vs_N$. Dans le modèle ZLF, $\vs_F$ a une composante
    transverse qui décroît sur $\lambda_J=\sqrt{2\hbar D_0/J_{sd}}$. Écrire le
    terme de mélange avec $\vs_N-\vs_F$ met alors en série l'absorption de
    l'interface et celle du volume de F2, et compte deux fois l'absorption
    si $G_{\uparrow\downarrow}$ la décrit déjà entièrement (valeur \emph{ab
    initio}). FeeLLGood suit par défaut la théorie des circuits (section
    \ref{sec:transverse}) ; l'écriture en série reste possible, avec un
    $G_r$ qui ne représente que la partie interfaciale (Petitjean
    \emph{et al.}).
\end{itemize}

\section{Le nœud espaceur : d'où vient la dépendance en \texorpdfstring{$\vm_0$}{m0}}

\subsection{Conservation du spin dans l'espaceur}

On suppose l'espaceur assez fin pour que $\vs_N$ y soit uniforme dans
l'épaisseur. Il est en revanche épais devant la maille latérale, donc on
écrit la relation en chaque point de $\Gamma$. On néglige la relaxation
dans N. Le spin qui entre dans N depuis F1 ressort vers F2 :
\begin{equation}
  \underbrace{P_1\frac{\mu_B}{e}J\,\vp + \mG_1(\vp)\,(\vs_{F1}-\vs_N)}_{\text{F1}\,\to\,\text{N}}
  \;=\;
  \underbrace{P_2\frac{\mu_B}{e}J\,\vm_0 + \mG_2(\vm_0)\,(\vs_N-\vs_F)}_{\text{N}\,\to\,\text{F2}}.
  \label{eq:noeud}
\end{equation}
Ici $J>0$ correspond à des électrons qui vont de F1 vers F2. Les indices 1
et 2 désignent les interfaces F1/N et N/F2, et $\vs_{F1}$ est
l'accumulation dans F1 à son interface. La section~\ref{sec:sF1} explique
comment l'obtenir.

On pose $\mG_1 = \mG_1(\vp)$, $\mG_2 = \mG_2(\vm_0)$ et
$\mA = \mG_1 + \mG_2$. L'équation \eqref{eq:noeud} est un système linéaire
$3\times3$ pour $\vs_N$ :
\begin{equation}
  \mA\,\vs_N = \frac{\mu_B}{e}J\left(P_1\vp - P_2\vm_0\right)
  + \mG_1\vs_{F1} + \mG_2\vs_F.
\end{equation}

\subsection{Élimination de \texorpdfstring{$\vs_N$}{sN} : une condition de Robin effective}

On reporte $\vs_N$ dans \eqref{eq:robin}. On utilise les identités
$\Id-\mG_2\mA^{-1} = \mG_1\mA^{-1}$ et
$\mG_2\mA^{-1}\mG_1 = \mG_1\mA^{-1}\mG_2$, et on obtient :
\begin{equation}
  \boxed{\;
  \vq = \frac{\mu_B}{e}\,J\,\bm{P}_{\!\mathrm{eff}}(\vm_0)
  + \mG_s\,\left(\vs_{F1}-\vs_F\right)
  \;}
  \label{eq:robin-eff}
\end{equation}
avec :
\begin{align}
  \bm{P}_{\!\mathrm{eff}}(\vm_0) &= \mG_2\mA^{-1}\,P_1\,\vp
  + \mG_1\mA^{-1}\,P_2\,\vm_0,
  \label{eq:Peff}\\
  \mG_s &= \mG_2\,\mA^{-1}\,\mG_1
  \quad\text{(conductance série des deux interfaces).}
\end{align}

C'est la réponse directe à la question de départ. La polarisation injectée
dans F2 n'est plus le vecteur fixe $P\vp$ : c'est un \textbf{vecteur
$\bm{P}_{\!\mathrm{eff}}$ qui dépend de $\vm_0$}, à la fois en direction et
en amplitude. Il interpole entre les deux polariseurs, avec un poids
donné par les conductances.

\subsection{Cas limites}

\begin{itemize}
  \item \textbf{Interface F1/N très résistive}
    ($\mG_1\ll\mG_2$) : $\mA\approx\mG_2$, donc
    $\bm{P}_{\!\mathrm{eff}}\to P_1\vp$ et $\mG_s\approx\mG_1\to0$.
    On retrouve exactement le modèle de départ \eqref{eq:neumann}.
    Celui-ci est donc la limite d'un polariseur résistif.
  \item \textbf{Interface N/F2 très résistive}
    ($\mG_2\ll\mG_1$) : $\bm{P}_{\!\mathrm{eff}}\to P_2\vm_0$. Le
    filtrage par F2 domine et le couple devient faible.
  \item \textbf{Configuration colinéaire} ($\vm_0=\pm\vp$) : toutes les
    matrices sont diagonales dans la même base, et on retrouve la formule
    scalaire de deux résistances en série (Valet--Fert).
  \item \textbf{F2 mince et absorption transverse totale} : la moyenne du
    couple sur l'épaisseur redonne la forme $\varepsilon(\theta)$ de
    l'équation~\eqref{eq:eps}, avec un $\Lambda$ fonction des conductances
    (Slonczewski 2002).
\end{itemize}

\section{Accès à l'accumulation dans F1}
\label{sec:sF1}

Les formules \eqref{eq:robin-eff}--\eqref{eq:Peff} font intervenir
$\vs_{F1}$, l'accumulation de spin dans F1 juste à son interface avec
l'espaceur. Il y a trois façons de l'obtenir, de la plus simple à la plus
complète.

\subsection{Option 1 : poser \texorpdfstring{$\vs_{F1}=0$}{sF1 = 0}}

C'est une approximation. Elle n'est valable que si la conductance de spin du
volume de F1, notée $g_F$ et définie plus bas, est grande devant la
conductance de l'interface : $g_F\gg G_{\parallel1}$. En général,
l'accumulation à l'interface n'est pas nulle : comme chez Valet--Fert, elle
décroît dans F1 sur la longueur $\lambda_{sf}$.

\subsection{Option 2 : élimination analytique (recommandée si F1 est fixe et uniforme)}

\paragraph{Pourquoi une seule inconnue.} L'aimantation de F1 est uniforme
et fixe : $\vm = \vp$ dans tout F1. Décomposons $\vs$ en une partie
longitudinale $s\,\vp$ et une partie transverse $\vs_\perp$
($\vs_\perp\cdot\vp = 0$). Dans \eqref{eq:zlf}, la partie transverse
précesse autour de $\vp$ (terme $\vs\times\vm/\tau_J$) et se déphase : elle
s'amortit sur la longueur $\lambda_J = \sqrt{2\hbar D_0/J_{sd}}\sim1$\,nm
à partir de l'interface F1/N. Le courant transverse qui arrive de
l'espaceur est donc absorbé par F1 en quelques couches atomiques, et c'est
la conductance de mélange $G_{r1}$, $G_{i1}$ de l'interface F1/N qui en rend
compte. Au-delà de $\lambda_J$, il ne reste que la partie longitudinale :
\begin{equation}
  \vs_{F1} = s_{F1}\,\vp,
\end{equation}
avec $s_{F1}$ la valeur de $s$ à l'interface. Il ne reste qu'\textbf{une
inconnue scalaire}, et le problème dans F1 est une diffusion colinéaire.
On suppose en outre F1 uniforme dans son plan : le problème est 1D.

\paragraph{Géométrie.} On place F1 en $-t_1<x<0$, l'interface F1/N en
$x=0$ et l'axe $\hat{\bm{x}}$ orienté de F1 vers F2. La face $x=-t_1$ est
en contact avec un réservoir (un métal épais, ou l'électrode), où
l'accumulation de spin est nulle : $s(-t_1)=0$. On note $\beta_1$, $\beta'_1$
les polarisations de volume de F1.

\paragraph{Courant de spin longitudinal.} Avec $\vm=\vp$ uniforme et
$\vs = s(x)\,\vp$, on a $\partial_i\vs = \vp\,\partial_i s$ et
$\vm\cdot\partial_i\vs = \partial_i s$. Le courant de spin
\eqref{eq:courant} devient
\begin{equation}
  Q_{i\alpha} = p_\alpha\left[-\beta_1\frac{\mu_B}{|e|}\,j_{e,i}
    - 2D_0\,(1-\beta_1\beta'_1)\,\partial_i s\right] :
\end{equation}
le spin transporté est partout parallèle à $\vp$. Sa composante utile est
le flux, selon $+\hat{\bm{x}}$, du spin projeté sur $\vp$,
$j_s(x) = Q_{x\alpha}\,p_\alpha$. Avec $J = -\bm{j}_e\cdot\hat{\bm{x}}$
(positif pour des électrons qui vont de F1 vers F2) et en notant $e=|e|$
comme dans le reste du document :
\begin{equation}
  j_s(x) = \beta_1\frac{\mu_B}{e}J - 2D_1\,s'(x),
  \qquad D_1 = D_0(1-\beta_1\beta'_1).
  \label{eq:F1-1D}
\end{equation}
Le premier terme est le spin porté par le courant de charge (dérive) ; il
est uniforme, car $J$ ne dépend pas de $x$ (conservation de la charge en
1D). Le second est la diffusion, avec le coefficient \emph{longitudinal}
$2D_1 = 2D_0(1-\beta_1\beta'_1)$, et non $2D_0$ : c'est lui qui fixe la
conductance de spin du volume de F1.

\paragraph{Équation de conservation.} En régime stationnaire,
\eqref{eq:zlf} projetée sur $\vp$ donne
$\partial_x j_s = -s/\tau_{sf}$, le terme de précession étant nul
($\vs\times\vm = s\,\vp\times\vp = \bm{0}$). Le terme de dérive étant
uniforme, il ne reste que la diffusion :
\begin{equation}
  2D_1\,s'' = \frac{s}{\tau_{sf}}
  \qquad\Longleftrightarrow\qquad
  s'' = \frac{s}{\lambda_1^2},
  \qquad \lambda_1^2 = 2D_1\,\tau_{sf}.
\end{equation}
$\lambda_1$ est la longueur de diffusion de spin \emph{longitudinale} de F1,
celle que l'on mesure (expériences de CPP-GMR, par exemple) : c'est elle,
et non $\sqrt{2D_0\tau_{sf}}$, qu'il faut donner comme paramètre.

\paragraph{Résolution.} La solution générale est
$s(x) = A\sinh\big((x+t_1)/\lambda_1\big) + B\cosh\big((x+t_1)/\lambda_1\big)$.
La condition de réservoir $s(-t_1)=0$ impose $B=0$, et la valeur à
l'interface $s(0) = s_{F1}$ donne $A = s_{F1}/\sinh(t_1/\lambda_1)$ :
\begin{equation}
  s(x) = s_{F1}\,\frac{\sinh\!\big((x+t_1)/\lambda_1\big)}{\sinh(t_1/\lambda_1)}.
\end{equation}
Sa dérivée à l'interface vaut
\begin{equation}
  s'(0^-) = \frac{s_{F1}}{\lambda_1}\,
  \frac{\cosh(t_1/\lambda_1)}{\sinh(t_1/\lambda_1)}
  = \frac{s_{F1}}{\lambda_1}\coth\frac{t_1}{\lambda_1}.
\end{equation}
En la reportant dans \eqref{eq:F1-1D}, le courant de spin qui sort de F1 par
l'interface vaut :
\begin{equation}
  j_s(0^-) = \beta_1\frac{\mu_B}{e}J - g_F\,s_{F1},
  \qquad
  g_F = \frac{2D_1}{\lambda_1}\coth\frac{t_1}{\lambda_1}.
  \label{eq:gF}
\end{equation}

\paragraph{Interprétation de $g_F$.} Le volume de F1 se comporte comme une
source de spin $\beta_1\frac{\mu_B}{e}J$ en parallèle avec une conductance
de spin $g_F$ (en m/s, dans les unités de $\vs$) : une accumulation
$s_{F1}>0$ à l'interface renvoie du spin vers le réservoir et diminue le
courant de spin qui sort de F1. Deux limites :
\begin{itemize}
  \item \textbf{F1 épais} ($t_1\gg\lambda_1$) : $\coth\to1$ et
    $g_F\to 2D_1/\lambda_1$. Seule une épaisseur $\lambda_1$ près de
    l'interface compte ; le réservoir est trop loin pour intervenir.
  \item \textbf{F1 mince} ($t_1\ll\lambda_1$) : $\coth(t_1/\lambda_1)\approx
    \lambda_1/t_1$ et $g_F\approx 2D_1/t_1$, très grand. Le réservoir,
    proche, court-circuite l'accumulation : $s_{F1}\to0$, et l'on retrouve
    l'option~1.
\end{itemize}
On écrit la continuité avec la composante longitudinale du courant à
l'interface F1/N :
\begin{equation}
  \beta_1\frac{\mu_B}{e}J - g_F\,s_{F1}
  = P_1\frac{\mu_B}{e}J + G_{\parallel1}\,\big(s_{F1} - \vs_N\cdot\vp\big),
\end{equation}
d'où
\begin{equation}
  s_{F1} = \frac{(\beta_1-P_1)\,\frac{\mu_B}{e}J + G_{\parallel1}\,\vs_N\cdot\vp}
  {g_F + G_{\parallel1}}.
\end{equation}
On remplace $s_{F1}$ dans le courant F1$\to$N. Le résultat a
\textbf{exactement la même forme qu'avec $\vs_{F1}=0$}, mais avec deux
paramètres renormalisés :
\begin{equation}
  \boxed{\;
  P_1^* = \frac{\beta_1\,G_{\parallel1} + P_1\,g_F}{g_F + G_{\parallel1}},
  \qquad
  G_{\parallel1}^* = \frac{g_F\,G_{\parallel1}}{g_F + G_{\parallel1}}
  \;}
  \label{eq:renorm}
\end{equation}
$G_{r1}$ et $G_{i1}$ ne changent pas.
\begin{itemize}
  \item $P_1^*$ est une moyenne de la polarisation de volume $\beta_1$ et de
    la polarisation d'interface $P_1$, pondérée par les conductances.
  \item $G_{\parallel1}^*$ est la conductance de l'interface et celle du
    volume mises en série.
\end{itemize}
En pratique, \textbf{on garde toutes les formules précédentes avec
$\vs_{F1}=0$, en remplaçant $P_1\to P_1^*$ et
$G_{\parallel1}\to G_{\parallel1}^*$} dans $\mG_1(\vp)$. On n'a rien à
simuler dans F1. Dans la limite $g_F\gg G_{\parallel1}$, on retrouve
$P_1^*\to P_1$, $G_{\parallel1}^*\to G_{\parallel1}$ et l'option~1. Le signe
du terme en $\beta_1$ dépend de la convention choisie pour le sens de $J$,
comme dans le reste du document.

\subsection{Paramètres de F1 dans FeeLLGood}
\label{sec:F1-feellgood}

FeeLLGood utilise le modèle ZLF avec $\beta = \beta' = P$, et écrit le flux
de diffusion $D_0\nabla\vs$, avec $D_0 = 2\sigma/(e^2N_0)$ (le facteur 2 du
modèle ZLF est inclus dans $D_0$). Dans les régions magnétiques :
\begin{itemize}
  \item la diffusion longitudinale vaut $D_0(1-P^2)$, la diffusion
    transverse $D_0$ (terme $-D_0P^2\,(\vm\cdot\partial_i\vs)
    (\vm\cdot\partial_i\bm{w})$ de la forme faible) ;
  \item la relaxation vaut $1/\tau_{sf} = D_0(1-P^2)/l_{sf}^2$, de sorte
    que $l_{sf}$ est la longueur de diffusion de spin longitudinale.
\end{itemize}
Ces deux points sont implémentés dans la branche \texttt{slonczewski}
(commit \texttt{ef95723}). Sur un cylindre aimanté uniformément, avec une
injection de spin selon $\vm$, la décroissance longitudinale mesurée vaut
$l_{sf}$ pour $P=0{,}3$ et $P=0{,}7$, et l'amplitude à l'entrée suit
$1/(1+P)$, comme le prévoit la théorie 1D.

Les paramètres de F1 sont $\sigma_1$, $P_{F1}$, $l_{sf,1}$ et $t_1$, et
les grandeurs de l'option~2 deviennent :
\begin{equation}
  \boxed{\;
  \beta_1 = P_{F1},\qquad
  \lambda_1 = l_{sf,1},\qquad
  g_F = \frac{2}{e^2N_0}\,
        \frac{\sigma_1\,(1-P_{F1}^2)}{l_{sf,1}}\,
        \coth\frac{t_1}{l_{sf,1}}
  \;}
  \label{eq:gF-feellgood}
\end{equation}
Ici $N_0$ est la densité d'états de l'\emph{échantillon} F2, et non celle
de F1 (voir la section~\ref{sec:conversion}) : la densité d'états de F1
n'intervient pas.

Il n'y a donc pas de paramètre $\beta_1$ séparé : c'est le $P$ de volume de
F1. Il ne faut pas le confondre avec $P_1$, la polarisation de
l'\emph{interface} F1/N ($\gamma_1$ en CPP-GMR), qui reste un paramètre à
part.

\subsection{Option 3 : simuler F1}

Cette option est nécessaire si F1 n'est pas uniforme, si la géométrie est
réellement 3D (courant non uniforme, F1 de taille latérale finie) ou si F1
est une structure plus complexe, par exemple un antiferromagnétique
synthétique.

On résout alors l'équation ZLF \eqref{eq:zlf} dans F1 comme dans F2.
$\vs_{F1}$ est la trace de cette solution sur l'interface F1/N. La relation
du nœud espaceur \eqref{eq:noeud} devient une condition qui relie les deux
domaines, et on résout les deux d'un seul bloc. Une variante consiste à
mailler tout l'empilement F1/N/F2 avec des paramètres qui changent selon la
couche (approche d'Abert \emph{et al.}).

\section{Mise en œuvre par éléments finis}

\subsection{Forme faible}

La diffusion de spin relaxe en $\tau_{sf}\sim\mathrm{ps}$, alors que
l'aimantation évolue en $\mathrm{ns}$. On résout donc la version
stationnaire de \eqref{eq:zlf} à $\vm$ fixé. Soit $\Omega$ le domaine F2
et $\vv$ une fonction test. On intègre par parties et on utilise
$\bm{n}\cdot Q = -\vq$ sur $\Gamma$. On suppose un flux nul sur
les bords latéraux et une condition de réservoir ($\vs=0$) ou de flux nul
à la sortie. On obtient :
\begin{multline}
  2D_0\int_\Omega\left[\partial_i s_\alpha
    - \beta\beta'\,m_\alpha\,(\vm\cdot\partial_i\vs)\right]\partial_i v_\alpha
  + \int_\Omega\left(\frac{\vs}{\tau_{sf}}
    + \frac{\vs\times\vm}{\tau_J}\right)\cdot\vv
  + \int_\Gamma\left(\mG_s\,\vs\right)\cdot\vv
  \\
  = -\int_\Omega\beta\frac{\mu_B}{|e|}\,j_{e,i}\,m_\alpha\,\partial_i v_\alpha
  + \int_\Gamma\left[\frac{\mu_B}{e}J\,\bm{P}_{\!\mathrm{eff}}(\vm_0)
    + \mG_s\,\vs_{F1}\right]\cdot\vv.
  \label{eq:faible}
\end{multline}
Avec l'option~2 de la section~\ref{sec:sF1}, on prend $\vs_{F1}=0$ et
le dernier terme disparaît. Par rapport au modèle de départ, il y a deux
changements :
\begin{enumerate}
  \item un \textbf{terme de masse de bord}
    $\int_\Gamma(\mG_s\vs)\cdot\vv$ dans la matrice ;
  \item un \textbf{second membre de bord} qui dépend de $\vm_0$, au lieu
    du flux constant $P\frac{\mu_B}{e}J\vp$.
\end{enumerate}
$\mG_s$ a une partie symétrique positive (termes $G_\parallel$ et $G_r$),
ce qui renforce la coercivité. Sa partie antisymétrique (terme $G_i$) ne
contribue pas à l'énergie.

\subsection{Couplage avec LLG}

À chaque pas de temps de LLG :
\begin{enumerate}
  \item lire $\vm$, et en particulier sa trace $\vm_0$ sur $\Gamma$ ;
  \item calculer en chaque point de $\Gamma$ les matrices
    $\mG_1(\vp)$ (avec $P_1^*$ et $G_{\parallel1}^*$ de
    \eqref{eq:renorm}, à faire une seule fois puisque $\vp$ est fixe),
    $\mG_2(\vm_0)$, $\mA^{-1}$, puis $\mG_s$ et
    $\bm{P}_{\!\mathrm{eff}}(\vm_0)$ (inversions $3\times3$ locales) ;
  \item réassembler les termes de bord de \eqref{eq:faible} et les termes
    de volume qui dépendent de $\vm$ ;
  \item résoudre le problème stationnaire pour $\vs$ ;
  \item calculer le couple $\bm{\tau}\propto\vm\times\vs$ et avancer LLG.
\end{enumerate}
Dans cette formulation, tout le courant $\vq$ entre dans F2, et le couple
vient entièrement du terme de volume $\vm\times\vs$. Pour décrire en plus
une absorption localisée à l'interface, il faut distinguer le courant côté N
du courant côté F2 : leur différence donne un couple surfacique (voir
Petitjean \emph{et al.}).

\subsection{Tests de validation}

\begin{itemize}
  \item Limite $\mG_1\ll\mG_2$ : retrouver les résultats du modèle de
    départ. \emph{Fait} (section~\ref{sec:validation}).
  \item F2 mince, $G_r$ grand : le couple moyenné sur l'épaisseur doit
    suivre $\varepsilon(\theta)\,\vm\times(\vm\times\vp)$, avec le
    $\Lambda$ correspondant aux conductances. \emph{Fait}
    (section~\ref{sec:validation}) : ce test fixe la conversion des
    conductances de mélange.
  \item Configuration colinéaire : comparer à la solution analytique de
    Valet--Fert. \emph{Fait} (section~\ref{sec:validation}).
  \item Renormalisation de F1 : simuler F1 explicitement (option~3) et
    vérifier que l'on retrouve le résultat obtenu avec $P_1^*$ et
    $G_{\parallel1}^*$ (option~2). \emph{Couvert indirectement} par le cas
    Valet--Fert, dont la solution analytique résout F1 explicitement.
\end{itemize}

\section{Implémentation dans FeeLLGood}
\label{sec:implementation}

L'option~2 est implémentée dans la branche \texttt{slonczewski} de
FeeLLGood, créée à partir de \texttt{normal\_current\_density}.

\subsection{Commits}

\begin{center}
\renewcommand{\arraystretch}{1.2}
\begin{tabular}{>{\ttfamily}l>{\raggedright\arraybackslash}p{11.5cm}}
  \toprule
  \textrm{Commit} & Contenu \\
  \midrule
  ef95723 & Diffusion anisotrope ($\beta=\beta'=P$) et relaxation
            $D_0(1-P^2)/l_{sf}^2$ dans les régions magnétiques. \\
  2ba96f0 & \texttt{src/slonczewskiInterface.h} : fonctions pures
            (matrice de conductance, renormalisation de F1, condition
            effective) et leurs tests unitaires. \\
  7c28c12 & Condition de Robin dans \texttt{spinAcc::solve()}
            (\texttt{stackBoundaryCondition}) et lecture du bloc
            \texttt{slonczewski} des surfaces. \\
  a7c409f & Cas de validation \texttt{validation/valet\_fert\_1d}. \\
  5887bbe & Interfaces données par $AR$ et $\gamma$, conversion faite par
            FeeLLGood avec la densité d'états de l'échantillon. \\
  3bc3af0 & Conductances de mélange converties en
            $2k\,G_{\uparrow\downarrow}$ (le facteur 2 manquait). \\
  d98d496 & Cas de validation angulaire
            \texttt{validation/slonczewski\_angle}. \\
  807fad5 & Signe de $\mathrm{Im}\,G_{\uparrow\downarrow}$ validé dans le cas
            angulaire. \\
  5a85673 & Continuité de $\vs$ entre régions de densités d'états
            différentes : documentation, avertissement en mode verbeux et
            cas de validation \texttt{validation/two\_regions\_1d}. \\
  cdc8f86 & Courant transverse à travers N/F2 : puits de la
            théorie des circuits par défaut, \texttt{transverse: series}
            en option. \\
  \bottomrule
\end{tabular}
\end{center}

\subsection{Réglages}

Le bloc \texttt{slonczewski} remplace \texttt{uP} sur la surface
d'injection. Il exige \texttt{jn} sur la même surface et y exclut
\texttt{uP} et \texttt{s} :
\begin{verbatim}
inlet:
  jn: 1e12            # jn > 0 : électrons de F1 vers F2
  slonczewski:
    p: [0, 0, 1]
    F1:  {sigma: 2e6, P: 0.5, l_sf: 20e-9, thickness: 15e-9}  # facultatif
    F1N: {AR: 0.5e-15, gamma: 0.75, Gr: 5e14, Gi: 0}  # AR en Ohm.m^2, Gr, Gi :
    NF2: {AR: 0.3e-15, gamma: 0.6,  Gr: 5e14,         # G_updown en S/m^2
          transverse: sink}                           # ou series (§ 11.3)
\end{verbatim}
\begin{itemize}
  \item \textbf{Avec \texttt{NF2}} (variante complète) : la condition
    dépend de $\vm_0$, et la surface doit être magnétique.
  \item \textbf{Sans \texttt{NF2}} (espaceur maillé, ou interface N/F2
    parfaite) : $\bm{P}_{\!\mathrm{eff}} = P_1^*\vp$ et
    $\mG_s = \mG_1(\vp)$, indépendants de l'aimantation.
  \item \textbf{Sans \texttt{F1}} : $\vs_{F1}=0$, sans renormalisation
    (option~1).
\end{itemize}

\subsection{Conversion des conductances}
\label{sec:conversion}

L'accumulation de spin de FeeLLGood (en A/m) est reliée à l'accumulation
de spin $\mu_s = (\mu_+-\mu_-)/2$ de Valet--Fert et de Brataas \emph{et
al.} par
\begin{equation}
  \vs = \mu_B\,N_0\,\frac{\mu_s}{2} = \mu_B\,N_0\,\frac{\mu_+-\mu_-}{4}.
\end{equation}
Avec cette relation, les équations de volume de Valet--Fert sont exactement
celles de FeeLLGood ($\beta=\beta'=P$, $D_0 = 2\sigma/(e^2N_0)$). Le
facteur $1/2$ vient de ce $D_0$, qui vaut deux fois le coefficient $D$ de
la relation d'Einstein $\sigma = e^2N_0D$ (c'est le \og $2D_0$ \fg{} du
modèle ZLF) : $\vs$ vaut la moitié de l'aimantation physique
$\mu_BN_0\mu_s$. Cela ne change aucun résultat, mais il faut s'en
souvenir pour interpréter \texttt{dens\_state}. Une
interface de produit résistance--surface $AR^*$ et d'asymétrie $\gamma$
($r_{\pm} = 2AR^*(1\mp\gamma)$) transporte le courant de spin
$\gamma\,J + [\mu_s]/(e\,AR^*)$ le long de son aimantation $\bm{u}$,
puisque $(1-\gamma^2)(G_\uparrow+G_\downarrow) = 1/AR^*$. En théorie des
circuits, le courant de spin transverse vaut
$2\,G_{\uparrow\downarrow}\,[\mu_s]_\perp/e$, où $G_{\uparrow\downarrow}$
est la conductance de mélange tabulée (par exemple
$\mathrm{Re}\,G_{\uparrow\downarrow}\approx0{,}5\cdot10^{15}$\,S/m² pour
Co/Cu). D'où les conductances en unités du problème de diffusion (m/s),
avec $k = 2/(e^2N_0)$ :
\begin{equation}
  \boxed{\;
  G_\parallel = \frac{k}{AR^*},\qquad
  P = \gamma,\qquad
  G_r + iG_i = 2k\,G_{\uparrow\downarrow}
  \;}
  \label{eq:conversion}
\end{equation}
Le facteur 2 de la partie transverse se vérifie sur une interface non
magnétique ($\gamma=0$, $G_\uparrow = G_\downarrow = G_{\uparrow\downarrow}
= G$) : le courant de spin y vaut $2G\,[\mu_s]/e$ dans toutes les
directions, donc $1/AR^* = 2G$, et la matrice de conductance doit être
isotrope, $G_\parallel = G_r$. C'est le test unitaire
\texttt{makeStack\_isotropic}.
$N_0$ est la densité d'états de la région de volume voisine de la surface,
c'est-à-dire de l'échantillon F2. L'empilement n'étant pas maillé, toutes
ses conductances (les deux interfaces et le volume de F1, éq.
\eqref{eq:gF-feellgood}) s'expriment d'abord en potentiel chimique, puis se
convertissent avec ce seul $N_0$. FeeLLGood fait cette conversion une seule
fois, à la construction du solveur (\texttt{spinAcc::makeStacks}) : il
retrouve le tétraèdre voisin de chaque triangle de la surface, et refuse
une surface qui borde des régions de densités d'états différentes.

\paragraph{Continuité de $\vs$ entre régions maillées.} FeeLLGood garde
$\vs$ continu entre les régions de volume, comme le courant de spin. Avec
$\vs = \mu_BN_0\,\mu_s/2$, c'est la continuité de $\mu_s$ de Valet--Fert
seulement si les régions ont la même \texttt{dens\_state} ; sinon, cela
impose à l'interface entre deux régions A et B
\begin{equation}
  N_{0,A}\,\mu_{s,A} = N_{0,B}\,\mu_{s,B},
  \label{eq:s-continu}
\end{equation}
sans résistance d'interface (le flux reste continu). C'est un choix de
modèle : pour retrouver la continuité de $\mu_s$ de Valet--Fert, on donne la
même \texttt{dens\_state} à toutes les régions, et chacune garde son propre
$\sigma$, donc son propre $D_0$. En mode verbeux, FeeLLGood signale les
régions voisines de \texttt{dens\_state} différentes
(\texttt{spinAcc::checkDensityOfStates}). Ce modèle est documenté dans la
classe \texttt{spinAcc}, dans \texttt{default-settings.yml} (clé
\texttt{dens\_state}) et dans \texttt{slonczewskiInterface.h}.

\paragraph{Courant transverse à travers N/F2.}
\label{sec:transverse}
La clé \texttt{transverse} de \texttt{NF2} choisit entre deux modèles :
\begin{itemize}
  \item \texttt{sink} (défaut), le \og puits \fg{} de la théorie des
    circuits : le courant transverse à travers N/F2 ne dépend que de
    $\vs_N$. $G_{\uparrow\downarrow}$ décrit toute l'absorption, et la
    diffusion dans F2 ne fait que répartir le couple sur $l_{sd}$. Dans
    l'élimination de l'espaceur, $\vs_F$ n'agit sur N/F2 que par sa partie
    longitudinale $\vm_0(\vm_0\cdot\vs_F)$ : la matrice de bord devient
    $\mG_s\,\vm_0\vm_0^{\mathsf T}$ (non symétrique ; le solveur BiCGStab
    l'accepte).
  \item \texttt{series} : le courant transverse dépend de
    $\vs_N-\vs_F$. L'interface ($1/G_{r2}$) et le volume de F2
    ($1/g_\perp$, avec $g_\perp = D_0\kappa$ et
    $\kappa^2 = (1-P^2)/l_{sf}^2 - i/l_{sd}^2$) sont deux résistances
    d'absorption en série ; $G_{\uparrow\downarrow}$ ne doit alors
    représenter que la partie interfaciale.
\end{itemize}
Dans le cas angulaire, $|2kG_{\uparrow\downarrow}| = 0{,}70\cdot10^6$\,m/s
et $|g_\perp| = 2{,}32\cdot10^6$\,m/s : le volume ajoute environ $30\,\%$
de résistance d'absorption, et le couple de type amortissement est de 15 à
$17\,\%$ plus faible avec \texttt{series} qu'avec \texttt{sink}.

\paragraph{Signe de la partie imaginaire.} La conductance de mélange est
celle de Brataas \emph{et al.} (Phys.\ Rep.\ \textbf{427}, 157, éq.~2) :
\begin{equation}
  G_{\uparrow\downarrow} = \frac{e^2}{h}\Big[M - \sum_{nm} r_\uparrow^{nm}
  \big(r_\downarrow^{nm}\big)^*\Big],
\end{equation}
où ↑ désigne le spin majoritaire, aligné sur $\vm$ (leur figure~2), comme le
canal majoritaire de FeeLLGood. Le courant de spin transverse de N vers F,
du côté N, vaut
$2\,\mathrm{Re}\,G_{\uparrow\downarrow}\,\vm\times(\bm{V}_S\times\vm)
+ 2\,\mathrm{Im}\,G_{\uparrow\downarrow}\,\bm{V}_S\times\vm$, avec
$\bm{V}_S = -\bm{\mu}_s/e$ puisque $\mu = -eV$. Dans les variables de
FeeLLGood, le terme imaginaire devient donc
$+2\,\mathrm{Im}\,G_{\uparrow\downarrow}\,
(\bm{\mu}_{s,F}-\bm{\mu}_{s,N})\times\vm/e$, et c'est le signe de
$\mG = \ldots - G_i[\vm]_\times$ dans \eqref{eq:G}. Le même signe se
retrouve à partir de la diffusion d'un seul canal : un spin transverse
$(|{\uparrow}\rangle+|{\downarrow}\rangle)/\sqrt2$ est réfléchi en
$(r_\uparrow, r_\downarrow)/\sqrt2$, et la polarisation transverse absorbée
vaut $1 - r_\uparrow^*r_\downarrow = \overline{G_{\uparrow\downarrow}}$. Le
cas angulaire le vérifie numériquement (section~\ref{sec:validation}).

\subsection{Validation}
\label{sec:validation}

\paragraph{Diffusion anisotrope.} Cylindre de 1000\,nm aimanté
uniformément, injection de spin selon $\vm$, $l_{sf}=10$\,nm : la
décroissance mesurée vaut $10{,}07$\,nm pour $P=0{,}3$ et $P=0{,}7$, et
l'amplitude à l'entrée suit $1/(1+P)$ (rapport $1{,}3077$ entre les deux
cas, exactement la valeur théorique).

\paragraph{Non-régression.} Avec $P_1=1$ et des conductances quasi nulles,
la condition de Robin redonne l'accumulation de l'injection \texttt{uP} à
$5\cdot10^{-8}$ près sur tout le maillage.

\paragraph{Valet--Fert 1D.} Le script
\texttt{validation/valet\_fert\_1d/valet\_fert.py} compare l'accumulation
de spin d'un fil F2 de $10\times10\times300$\,nm à la solution analytique
de Valet--Fert pour l'empilement F1 ($t_1=15$\,nm, $\beta_1=0{,}5$,
$\lambda_1=20$\,nm) / F1/N ($AR^*=0{,}5$\,f$\Omega$m², $\gamma=0{,}75$) /
espaceur d'épaisseur nulle / N/F2 ($0{,}3$\,f$\Omega$m², $\gamma=0{,}6$) /
F2 ($\beta=0{,}3$, $\lambda=10$\,nm), aimantations parallèles (P) ou
antiparallèles (AP) selon $z$. La solution analytique est écrite
indépendamment du code (système linéaire $5\times5$, par canal).
\begin{center}
\renewcommand{\arraystretch}{1.2}
\begin{tabular}{lccc}
  \toprule
  & \multicolumn{2}{c}{$s_z$ à l'entrée, FeeLLGood / Valet--Fert} & \\
  Maillage & P & AP & Erreur max \\
  \midrule
  $h=2$\,nm & $45{,}69$ / $45{,}86$\,A/m & $-163{,}99$ / $-164{,}60$\,A/m & $0{,}82\,\%$ \\
  $h=1$\,nm & $45{,}81$ / $45{,}86$\,A/m & $-164{,}43$ / $-164{,}60$\,A/m & $0{,}26\,\%$ \\
  \bottomrule
\end{tabular}
\end{center}
L'erreur décroît à peu près comme $h^2$ (éléments P1). Avec le facteur 1
au lieu de 2 dans \eqref{eq:conversion}, l'erreur monte à $8{,}5\,\%$ en
configuration parallèle : le test échoue, il fixe donc bien la conversion.
Il couvre aussi la renormalisation de F1, la mise en série des interfaces,
les signes de $j_n$ et des polarisations, la diffusion anisotrope et la
relaxation dans F2. Il ne couvre ni les conductances de mélange (sans effet
en colinéaire), ni des densités d'états différentes entre couches.

\paragraph{Cas angulaire sur F2 mince.} Dans le répertoire
\texttt{validation/slonczewski\_angle}, le script
\texttt{slonczewski\_angle.py} traite une
couche F2 de $2\times2\times3$\,nm maillée avec $h=0{,}2$\,nm (pour
résoudre $l_{sd}=1$\,nm), d'aimantation $\vm=\hat z$ et $P=0{,}5$, suivie
d'un réservoir. Le polariseur vaut $\vp = (\sin\theta, 0, \cos\theta)$,
avec deux interfaces de même $AR^*=0{,}5$\,f$\Omega$m², $\gamma=0{,}75$ et
$\mathrm{Re}\,G_{\uparrow\downarrow}=0{,}55\cdot10^{15}$\,S/m², et
$\mathrm{Im}\,G_{\uparrow\downarrow} = 0{,}15\cdot10^{15}$ (F1/N) et
$0{,}25\cdot10^{15}$\,S/m² (N/F2), valeurs plus grandes que d'habitude pour
que le signe compte. La référence est indépendante du code : les interfaces
sont écrites avec les formules de Brataas \emph{et al.} en $\mu_s$, terme
imaginaire compris, et les équations 1D de F2 sont résolues exactement par
modes propres ; elle suit les deux modèles du courant transverse à travers
N/F2. Le couple $\vm\times\vs$ est de type amortissement pour $s_y$
(direction $\vm\times(\vm\times\vp)$) et de type champ pour $s_x$. Avec
\texttt{transverse: sink} (défaut) :
\begin{center}
\renewcommand{\arraystretch}{1.2}
\begin{tabular}{rccc}
  \toprule
  $\theta$ & Amortissement $\int s_y\,dz$ & Champ $\int s_x\,dz$ & Erreur max \\
  & FeeLLGood / référence ($10^{-9}$\,A) & FeeLLGood / référence ($10^{-9}$\,A) & sur $\vs$ \\
  \midrule
  30° & $4{,}7941$ / $4{,}7930$ & $1{,}5079$ / $1{,}5266$ & $0{,}46\,\%$ \\
  60° & $7{,}9075$ / $7{,}9065$ & $2{,}2985$ / $2{,}3291$ & $1{,}1\,\%$ \\
  90° & $8{,}1529$ / $8{,}1531$ & $2{,}0235$ / $2{,}0547$ & $0{,}52\,\%$ \\
  120° & $5{,}9018$ / $5{,}9031$ & $1{,}1078$ / $1{,}1300$ & $0{,}23\,\%$ \\
  150° & $2{,}8398$ / $2{,}8410$ & $0{,}3526$ / $0{,}3631$ & $0{,}10\,\%$ \\
  \bottomrule
\end{tabular}
\end{center}
Avec \texttt{transverse: series}, les résultats sont ceux d'avant ce choix
(erreur maximale $1{,}1\,\%$), par exemple à 90° :
$6{,}8434$ / $6{,}8403\cdot10^{-9}$\,A pour le couple d'amortissement. Le
couple d'amortissement diffère de 15 à $17\,\%$ entre les deux modèles :
le test les distingue. L'accumulation transverse s'annule à 0 et 180°.
L'erreur décroît comme $h^2$ ($1{,}1\,\%$ avec $h=0{,}2$\,nm,
$0{,}27\,\%$ avec $h=0{,}1$\,nm, vérifié avec \texttt{series} et
$\mathrm{Im}\,G_{\uparrow\downarrow}=0$). Le test échoue aussi si la
conversion ou le signe est faux :
\begin{itemize}
  \item avec l'ancienne conversion $k\,G_{\uparrow\downarrow}$ au lieu de
    $2k\,G_{\uparrow\downarrow}$ ($\mathrm{Im}\,G_{\uparrow\downarrow}=0$),
    l'erreur atteint $170\,\%$ ;
  \item avec le signe opposé du terme en $\mathrm{Im}\,G_{\uparrow\downarrow}$
    dans la référence, l'erreur atteint $13\,\%$, et le couple de champ
    diffère d'un facteur 2 à 150°.
\end{itemize}
Le couple de type amortissement suit la forme de Slonczewski
$\varepsilon(\theta)\sin\theta$ à $5\,\%$ près, avec
$\Lambda^2\approx0{,}53$ (\texttt{sink}) ou $0{,}54$ (\texttt{series}) : le
couple est plus fort près de la configuration parallèle pour cette
structure asymétrique.

\paragraph{Deux régions de densités d'états différentes.} Le script
\texttt{two\_regions.py}, dans le répertoire
\texttt{validation/}\allowbreak\texttt{two\_regions\_1d}, traite un fil de
$4\times4\times300$\,nm ($h=1$\,nm) : un espaceur N non magnétique maillé
($0<z<40$\,nm, $\sigma=6\cdot10^7$\,S/m, \texttt{dens\_state}
$0{,}6\cdot10^{47}$\,J$^{-1}$m$^{-3}$, $l_{sf}=50$\,nm), puis F2
($\sigma=4\cdot10^6$\,S/m, \texttt{dens\_state} $1{,}344\cdot10^{47}$,
$P=0{,}5$, $l_{sf}=10$\,nm), puis un réservoir. Le spin est injecté en
$z=0$ par un empilement sans interface N/F2 (l'espaceur étant maillé) :
F1 réservoir / F1/N ($AR^*=0{,}5$\,f$\Omega$m², $\gamma=0{,}75$),
aimantations parallèles (P) ou antiparallèles (AP) selon $z$. La référence
est écrite dans les variables de Valet--Fert avec la règle d'interface
\eqref{eq:s-continu}.
\begin{center}
\renewcommand{\arraystretch}{1.2}
\begin{tabular}{lccc}
  \toprule
  & \multicolumn{2}{c}{$\vs$ continu (FeeLLGood)} & $\mu_s$ continu \\
  & $s_z$ à l'interface, FeeLLGood / modèle & Erreur max & $s_N$ / $s_F$ à l'interface \\
  \midrule
  P  & $-7{,}130$ / $-7{,}131$\,A/m & $0{,}14\,\%$ & $-6{,}40$ / $-14{,}34$\,A/m \\
  AP & $-23{,}365$ / $-23{,}369$\,A/m & $0{,}14\,\%$ & $-20{,}98$ / $-47{,}00$\,A/m \\
  \bottomrule
\end{tabular}
\end{center}
FeeLLGood suit le modèle à $\vs$ continu à $0{,}14\,\%$ près. Avec des
\texttt{dens\_state} égales, les deux modèles coïncident ($10^{-15}$). Ici
($N_{0,F}/N_{0,N}=2{,}24$), ils diffèrent de $100\,\%$ de $\max|\vs|$, du
côté F2 de l'interface : avec $\mu_s$ continu, $\vs$ y saute d'un facteur
$2{,}24$, et l'accumulation dans F2 à l'interface est environ deux fois
celle du modèle de FeeLLGood. Le couple dans F2 en dépend directement.

\subsection{Points ouverts}

\begin{enumerate}
  \item \textbf{Résolution transverse} : l'absorption transverse dans F2 se
    fait sur $l_{sd}\sim1$\,nm ; un maillage plus grossier ne la résout pas.
\end{enumerate}

\section*{Synthèse}

\begin{center}
\renewcommand{\arraystretch}{1.3}
\begin{tabular}{>{\raggedright}p{3.2cm}>{\raggedright}p{4.8cm}>{\raggedright\arraybackslash}p{5.6cm}}
  \toprule
  & \textbf{Modèle de départ} & \textbf{Modèle généralisé} \\
  \midrule
  Condition sur $\Gamma$ & Neumann : $\vq = P\frac{\mu_B}{e}J\vp$
    & Robin : $\vq = \frac{\mu_B}{e}J\bm{P}_{\!\mathrm{eff}}(\vm_0)
      + \mG_s(\vs_{F1}-\vs_F)$ \\
  Rôle de $\vp$ & direction du flux & source de $\vs_N$ (nœud espaceur) \\
  Dépendance angulaire & aucune
    & apparaît d'elle-même (retour de spin, réflexions) \\
  Paramètres & $P$
    & $AR^*$, $\gamma$, $G_r$, $G_i$ pour F1/N et N/F2 ;
      $\sigma_1$, $P_{F1}$, $l_{sf,1}$, $t_1$ pour F1
      (sections~\ref{sec:F1-feellgood} et~\ref{sec:conversion}) \\
  \bottomrule
\end{tabular}
\end{center}

\section*{Références}

\begin{itemize}
  \item J.\,C.~Slonczewski, \emph{Current-driven excitation of magnetic
    multilayers}, J.~Magn.\ Magn.\ Mater.\ \textbf{159}, L1 (1996).
  \item J.\,C.~Slonczewski, \emph{Currents and torques in metallic magnetic
    multilayers}, J.~Magn.\ Magn.\ Mater.\ \textbf{247}, 324 (2002) :
    origine de la formule à paramètre $\Lambda$.
  \item T.~Valet, A.~Fert, Phys.\ Rev.\ B \textbf{48}, 7099 (1993) :
    magnétorésistance géante en géométrie CPP, modèle de diffusion
    colinéaire.
  \item S.~Zhang, P.\,M.~Levy, A.~Fert, Phys.\ Rev.\ Lett.\ \textbf{88},
    236601 (2002) : modèle de diffusion ZLF.
  \item A.~Brataas, G.\,E.\,W.~Bauer, P.\,J.~Kelly, Phys.\ Rep.\
    \textbf{427}, 157 (2006) : théorie des circuits magnéto-électroniques.
  \item C.~Petitjean, D.~Luc, X.~Waintal, Phys.\ Rev.\ Lett.\
    \textbf{109}, 117204 (2012) : courants transverses en dérive-diffusion
    et conditions d'interface cohérentes.
  \item C.~Abert \emph{et al.}, Sci.\ Rep.\ \textbf{5}, 14855 (2015) :
    diffusion de spin 3D dans tout l'empilement en micromagnétisme.
\end{itemize}

\end{document}
